% =============================================================================
% Capítulo 2 — Trabalhos Relacionados   (peso: médio)
% Fonte: plano §8 (integração de LLMs e frameworks; arquiteturas; lock-in fora
% de IA; medição energética; quadro com coluna de comparação controlada).
% -----------------------------------------------------------------------------
% Redigido em 02/11/2026 com as skills trabalhos-relacionados-tcc e
% fundamentacao-teorica-tcc. Fica entre a Introdução e a Fundamentação.
% Trechos que comparavam a proposta com outros trabalhos saíram do Cap. de
% Fundamentação e vieram para cá: Mo et al. (2023) como precedente; a camada
% de Kandogan et al. (2025); o posicionamento diante do catálogo de Liu et al.
% (2025); o limite dos estudos de energia na pilha de inferência.
% -----------------------------------------------------------------------------
% Registro da busca (T09), 02/11/2026. Período: 2023 a 2026, exceto os
% trabalhos fora de IA (Kaur 2017, Feitosa 2017), escolhidos pelo tema.
%  ACM DL, pelo Portal CAPES (acesso UFC), campo Abstract, filtro 2023-2026;
%  "dl" = ACM Full-Text Collection, "guide" = ACM Guide (inclui outras editoras):
%   A1 ("large language model" OR LLM) AND ("vendor lock-in" OR "abstraction
%      layer" OR middleware OR gateway OR "provider-agnostic" OR
%      "model-agnostic") AND (architecture OR integration) ............ dl 61
%   A2 ("large language model" OR LLM OR "language model") AND ("vendor
%      lock-in" OR "lock-in" OR "provider-agnostic" OR "switching provider" OR
%      "provider switching" OR "model migration" OR portability OR
%      replaceability OR substitutability) ............................ dl 439
%   A3 Title:(LLM OR "language model(s)" OR "foundation model") AND
%      Title:(gateway OR middleware OR abstraction OR "lock-in" OR portability
%      OR "provider-agnostic" OR migration) ................ dl 82, guide 194
%   A4 ("large language model" OR LLM) AND ("software architecture" OR
%      "architectural decision" OR "reference architecture" OR "architectural
%      tactic") AND (integrated OR integration) ............. dl 10, guide 20
%   A5 energy AND ("agentic framework(s)" OR "agent framework" OR LangChain OR
%      AutoGen OR "orchestration framework" OR "LLM framework") .. dl 7, guide 26
%   A6 energy AND ("design pattern(s)" OR "architectural tactic(s)" OR
%      "abstraction layer" OR indirection) AND (consumption OR overhead OR
%      efficiency) ......................................... dl 33, guide 118
%   A7 ("large language model" OR LLM) AND (LangChain OR LiteLLM OR "LLM
%      framework(s)" OR "LLM library" OR "LLM SDK" OR "LLM API") AND (empirical
%      OR developers OR maintenance OR evolution) ......... dl 52, guide 103
%  IEEE Xplore, pelo Portal CAPES, campo Abstract, 2023-2026:
%   I1 = A1 ............................................................ 90
%   I2 (energy OR "power consumption") AND ("large language model" OR LLM OR
%      agentic) AND (framework OR orchestration OR middleware OR "abstraction
%      layer" OR "software architecture") AND (overhead OR "energy
%      consumption") ................................................... 171
%  Complementares: OpenAlex e Crossref (APIs), busca na web, SBC OpenLib por
%  evento (SBES, SBCARS, SEMISH, WSCAD/SSCAD, WCAMA, SBSI; a busca geral da
%  SBC exige login). Triagem por título e resumo; inclusão: revisado por pares
%  e com a mesma motivação, objeto ou método da proposta.
% Metadados dos trabalhos novos conferidos no Crossref em 02/11/2026.
% Lidos: resumos de todos os trabalhos novos (Feitosa et al. e Mo et al. pelo
% Portal CAPES); texto da versão arXiv de Bucaioni et al. (2501.12904v1).
% Atualização de 02/11/2026 (tasks T01, T06, T09 da revisão v1):
%  - lidos na versão arXiv: Kandogan et al. (2504.08148v1: sem comparação com
%    integração direta, sem esforço de troca, sem custo medido); Tripathy et
%    al. (2512.09543: razão 9,4x e taxas de resolução de 4% e 0%); Xiao et al.
%    (2501.14402: resultados e discussão); Lin et al. (2604.21083: dez gateways);
%  - lido só o resumo (OpenAlex): Matinpour et al. (números de latência e
%    energia conferidos no resumo); Morais et al. (negativas reformuladas como
%    "o resumo não relata"). Os dois estão no Portal CAPES (ACM DL): ler o texto
%    integral antes de confirmar ou desfazer a reformulação.
% NÃO lidos na íntegra: Morais et al. e Matinpour et al.
% Triados e não usados: preprints sem venue (arXiv 2609.31288, "When the
% Model Retires"; arXiv 2606.10702, "Watts and Debts", registered report);
% ArcAssist-Web (WebMedia 2026, IA no apoio à decisão arquitetural);
% Connolly Bree e Ó Cinnéide (JSEP 2025, SLR redundante com Poy et al.);
% Barbarroxa et al. (PAAMS 2024, compara frameworks só pela qualidade da
% saída); demais resultados fora do tema (aplicações de LLM, segurança,
% inferência em hardware).
% =============================================================================
\chapter{Trabalhos Relacionados}
\label{cap:trabalhos-relacionados}

Este capítulo compara a proposta com os trabalhos que compartilham a sua
motivação, o seu objeto ou o seu método. Os trabalhos foram identificados em
buscas na ACM Digital Library e no IEEE Xplore, acessados pelo Portal de
Periódicos CAPES, complementadas pelas bases abertas OpenAlex e Crossref e por
buscas por evento na Biblioteca Digital da Sociedade Brasileira de Computação
(SBC OpenLib), em novembro de 2026. As buscas cobriram publicações de 2023 a
2026 e combinaram termos de três grupos: camadas de integração com modelos de
linguagem (\textit{abstraction layer}, \textit{middleware}, \textit{gateway},
\textit{provider-agnostic}, \textit{vendor lock-in}), arquitetura de
\textit{software} de sistemas integrados a LLMs, e custo energético de
\textit{frameworks}, padrões e táticas arquiteturais (decisões de projeto
tomadas para influenciar um atributo de qualidade, como o consumo de energia
\cite{bass2021}). As \textit{strings}, a
data, o número de resultados e os critérios de triagem estão no
Apêndice~\ref{ap:busca}. Consideraram-se apenas
publicações revisadas por pares; dois trabalhos anteriores a esse período,
fora do domínio da IA, foram incluídos pelo tema.

Os trabalhos se organizam em quatro grupos, um para cada parte da pergunta de
pesquisa: as arquiteturas e camadas propostas para integrar modelos a sistemas
de \textit{software} (Seção~\ref{sec:tr-arquiteturas}); as camadas de
abstração usadas contra a dependência de fornecedor fora do domínio da IA
(Seção~\ref{sec:tr-lockin}); o custo energético de decisões arquiteturais
(Seção~\ref{sec:tr-energia-arquitetura}); e o consumo de energia em sistemas
com modelos de linguagem e agentes (Seção~\ref{sec:tr-energia-llm}). A
Seção~\ref{sec:tr-sintese} reúne as diferenças em um quadro comparativo.

% -----------------------------------------------------------------------------
\section{Arquiteturas e camadas para integrar modelos de linguagem}
\label{sec:tr-arquiteturas}

O trabalho mais próximo do projeto da camada é o catálogo de padrões
arquiteturais para agentes baseados em modelos de fundação de
\citeonline{liu2025}, que reúne 18 padrões identificados por revisão
sistemática da literatura e descreve, para cada um, contexto, forças e
consequências. A relação com a presente proposta está nas responsabilidades da
camada: parte delas tem correspondência, completa ou por analogia, com padrões
do catálogo, que serve de referência para delimitá-las. De forma semelhante,
este trabalho trata a integração com o modelo como problema de arquitetura e
discute os seus compromissos. Diferencia-se em dois pontos. O catálogo avalia
os compromissos de forma qualitativa, e os próprios autores observam que a
maior parte dos benefícios e desvantagens listados são atributos de qualidade
que ainda exigem quantificação; este trabalho mede dois deles. Além disso, o
objeto do catálogo é o projeto interno do agente, e não a camada que liga a
aplicação ao provedor do modelo.

\citeonline{bucaioni2025} propõem, em caráter preliminar, uma arquitetura de
referência funcional para sistemas integrados a LLMs, organizada em quatro
camadas: apresentação, lógica de aplicação, integração com o LLM e
gerenciamento de dados. A arquitetura é confrontada com preocupações como
privacidade, segurança, modularidade e interoperabilidade e é verificada por
correspondência com três sistemas de código aberto. Na mesma linha,
\citeonline{bandeira2026} propõem uma arquitetura de referência para
aplicações concebidas desde o início com o LLM como núcleo de orquestração,
ilustrada por um assistente para revisões de literatura. Esses trabalhos se
relacionam com a presente proposta por isolar a integração com o modelo em uma
camada ou núcleo próprio, separado da lógica de aplicação, que é a decisão
avaliada aqui. A diferença está no tipo de avaliação e no que a camada
resolve: as duas propostas são verificadas por mapeamento ou ilustração, sem
medição, e a troca de provedor não é o problema que endereçam; neste trabalho,
a troca de provedor é o benefício medido, e o custo da camada também é medido.

Mais próximo quanto ao objetivo está o trabalho de \citeonline{morais2026}, que
propõe especificar agentes de IA como artefatos declarativos, independentes da
tecnologia, do modelo e do provedor, e separar a declaração, o despacho e a
execução do agente, de modo que a integração saia do código imperativo de
orquestração. O modelo é materializado em uma arquitetura de execução com
registro de agentes, contratos, telemetria e substituição transparente do
modelo, e um estudo de viabilidade com agentes de análise de
\textit{code smells} baseados em LLMs mostra que papéis diferentes
compartilham o mesmo fluxo de composição, preservando a portabilidade entre
configurações de modelo. A semelhança com a presente proposta é direta: ambos
interpõem uma camada para que a aplicação não dependa do provedor e tratam a
troca de modelo como requisito. Diferenciam-se pela avaliação, pois o estudo de
viabilidade de \citeonline{morais2026} demonstra que a substituição é possível,
mas o resumo do trabalho não relata medição do esforço de fazê-la, comparação
com uma integração direta nem medição do custo da camada.

No domínio dos agentes, \citeonline{kandogan2025} propõem uma arquitetura de
referência para IA composta em ambiente corporativo, na qual registros de
agentes e de dados descrevem os modelos, APIs e serviços disponíveis e abstraem
provedores e APIs proprietárias, enquanto fluxos de mensagens
(\textit{streams}) orquestram a execução. A arquitetura é ilustrada por um caso
de uso. A semelhança com a presente proposta está no papel de intermediário: a
aplicação se relaciona com o registro, e não com cada provedor, como aqui se
relaciona com o \harness{}. A diferença é que a proposta de
\citeonline{kandogan2025} não é comparada a uma integração direta, de modo que
o efeito da abstração sobre o esforço de troca ou sobre o custo de execução não
é conhecido: o caso de uso serve para demonstrar a arquitetura, e os próprios
autores reconhecem que muito ainda está por projetar.

A mesma ideia aparece, na prática, sob a forma de \textit{gateways}: serviços
que oferecem uma API única para modelos de vários provedores.
\citeonline{sheneman2026} descrevem um \textit{gateway} de código aberto que
traduz entre os protocolos de cliente da OpenAI, do Ollama e da Anthropic e
distribui as requisições entre GPUs de uma universidade, e
\citeonline{rahaman2026} descrevem um \textit{gateway} institucional montado
sobre ferramentas abertas como o LiteLLM e o vLLM, compatível com a API da
OpenAI, com testes de carga preliminares. Os dois relatos confirmam que a
camada entre aplicação e provedor é adotada na prática, mas descrevem a
infraestrutura e a implantação, sem compará-la com o acesso direto.
\citeonline{lin2026}, por sua vez, auditaram dez \textit{gateways} comerciais
de terceiros e encontraram substituições silenciosas de modelo, divergências
de cobrança e instabilidade de latência, o que mostra o outro lado da
intermediação: a aplicação passa a depender também do comportamento da camada.
A presente proposta difere desses trabalhos por ser uma camada dentro da
própria aplicação, sob controle de quem a desenvolve, e por medir, e não apenas
descrever, o que ela traz e o que ela custa.

Em vez de construir a camada, uma equipe pode adotar um \textit{framework} de
orquestração. \citeonline{zhu2026} comparam seis desses \textit{frameworks},
entre eles LangGraph, CrewAI, AutoGen e o OpenAI Agents SDK, e os protocolos MCP
e A2A, discutindo gerenciamento de estado, tratamento de erros e segurança; ao
fazê-lo, distinguem os achados relatados pelos fornecedores dos sustentados
por avaliação revisada por pares. O \textit{survey} se relaciona com este
trabalho porque os \textit{frameworks} são a alternativa pronta a uma camada
própria, e a distinção entre os dois tipos de evidência reforça a motivação
para medir. Este trabalho, contudo, não compara \textit{frameworks}: implementa
uma camada própria para que a única diferença entre as duas versões da
aplicação seja a presença dela, o que um \textit{framework} completo, com
decisões de projeto próprias, não permitiria isolar.

% -----------------------------------------------------------------------------
\section{Camadas de abstração contra a dependência de fornecedor}
\label{sec:tr-lockin}

A dependência de fornecedor é tratada há mais tempo em computação em nuvem.
A revisão de \citeonline{kaur2017}, apresentada na
Seção~\ref{sec:arquitetura}, é o ponto de partida do tema. Ela difere da presente
proposta no domínio, serviços de nuvem e não provedores de modelos, e no tipo
de estudo, uma revisão, e não um experimento.

O precedente empírico mais próximo é o de \citeonline{mo2023}. Os autores
investigam bibliotecas de abstração de serviços de nuvem no acesso ao
armazenamento de objetos a partir de funções \textit{serverless}: comparam o
desempenho de sete funções na AWS e no Google Cloud escritas sobre a biblioteca
Apache jclouds e sobre as APIs específicas de cada plataforma, avaliam
métricas de qualidade do código e conduzem um estudo com 42 estudantes, que
migram código entre plataformas, para verificar o efeito da biblioteca sobre a
portabilidade, medida pelo tempo, pelo esforço percebido e pelo sucesso da
tarefa. O desenho é o mesmo deste trabalho: duas versões de uma aplicação, com
e sem a camada, comparadas quanto ao benefício da abstração e ao preço que ela
cobra. Também é comum o uso de métricas de código, que lá avaliam a qualidade
do código escrito com e sem a biblioteca. As diferenças estão no domínio, no
custo medido e em duas escolhas de método. Os provedores de modelos diferem não
só na API de acesso, mas no formato das mensagens e do ciclo de ferramentas,
que se repete várias vezes por requisição agêntica, isto é, a que envolve o
ciclo de ferramentas (Seção~\ref{sec:llm-agentes}); e o preço da camada é medido aqui em energia, e
não em desempenho.

A primeira escolha de método é a medida do esforço de troca. Em
\citeonline{mo2023}, ele foi medido com pessoas; aqui, por métricas de código
aplicadas pelo próprio autor, sem participantes. O trabalho perde a validade de
construto do esforço (o tempo e o esforço percebido de quem realmente migra o
código) e ganha um procedimento que dispensa recrutar participantes e que se
repete sem a variação entre pessoas. A segunda é a abstração avaliada: lá, uma
biblioteca de terceiros pronta; aqui, uma camada construída pelo autor. O
trabalho perde a possibilidade de atribuir o resultado a uma biblioteca
consolidada e ganha o isolamento da camada, pelo mesmo argumento dado a
\citeonline{zhu2026}: as duas versões da aplicação diferem apenas pela presença
dela.

% -----------------------------------------------------------------------------
\section{Custo energético de decisões arquiteturais}
\label{sec:tr-energia-arquitetura}

A hipótese de que uma decisão de projeto tenha custo energético mensurável já
foi estudada fora da IA. No mapeamento sistemático de \citeonline{poy2025}
sobre o efeito de padrões de projeto, \textit{code smells} e refatorações no
consumo de energia, o uso de padrões de projeto aparece, de modo geral, com
impacto negativo, embora nem todos tenham efeito claro, o que, segundo os
autores, exige mais estudos. Um exemplo desse tipo de estudo é o de
\citeonline{feitosa2017}, que compararam o consumo de 169 métodos que
participam dos padrões \textit{State}/\textit{Strategy} e \textit{Template
Method}, em dois projetos de código aberto, com soluções alternativas
funcionalmente equivalentes, obtidas por refatoração. A alternativa sem o
padrão consumiu menos energia na maior parte dos casos, mas, quando o método é
grande ou invoca muitos outros, a solução com o padrão teve consumo semelhante
ou menor. Esses trabalhos se relacionam com a presente proposta por tratar a
indireção introduzida por uma estrutura de projeto como possível fonte de
consumo e por compará-la com uma alternativa equivalente sem ela; o achado de
\citeonline{feitosa2017} sugere, além disso, que o peso relativo da indireção
cai quando o trabalho envolvido cresce. Por analogia, é de esperar que o mesmo
ocorra em uma requisição a um LLM, mas essa extensão é uma hipótese deste
trabalho, e não um achado dos autores. Diferenciam-se pela escala da decisão: padrões de projeto no nível do
código, e não uma camada arquitetural entre a aplicação e um serviço externo
cuja execução domina o custo da requisição.

No nível da arquitetura, \citeonline{xiao2025} avaliam, em experimento
controlado sobre o sistema de microsserviços Online Boutique, três táticas e
três padrões arquiteturais voltados a reduzir o consumo de energia, sob três
níveis de carga. Todas as técnicas reduziram o consumo em ao menos um nível de
carga, em até 5,6\%, mas três delas prejudicaram a manutenibilidade, uma a
favoreceu e duas não a alteraram. Esse é o trabalho mais próximo quanto à
forma da pergunta: confronta, para a mesma decisão arquitetural, energia e
manutenibilidade, que são as duas dimensões deste trabalho. Os autores
resumem os compromissos de cada técnica em uma tabela comparativa e
recomendam avaliar a carga do sistema antes de adotá-las, mas não propõem um
critério para decidir quando a economia de energia compensa a perda de
manutenibilidade; o critério de decisão do objetivo (e) não é, portanto,
herdado desse estudo, e sim contribuição própria deste trabalho. A direção do
compromisso, porém, é a inversa. Lá, as técnicas são adotadas para economizar
energia e podem custar manutenibilidade; aqui, a camada é adotada para
facilitar a troca de provedor e pode custar energia. Para sistemas com
aprendizado de máquina, \citeonline{jarvenpaa2024} reuniram 30 táticas
arquiteturais verdes, derivadas de 51 publicações e validadas com
especialistas; o catálogo, contudo, orienta o projeto para reduzir o consumo e
não avalia o efeito energético de uma decisão tomada por outro motivo, como a
de interpor uma camada de integração.

% -----------------------------------------------------------------------------
\section{Energia em sistemas com modelos de linguagem e agentes}
\label{sec:tr-energia-llm}

A medição de energia na inferência de LLMs concentra-se na pilha que executa o
modelo. \citeonline{vellaisamy2026} decompõem a sobrecarga de orquestração
visível no processador hospedeiro, separando-a do trabalho executado na GPU,
e mostram que ela é um custo oculto da latência da inferência.
\citeonline{ifath2026} caracterizam compromissos entre desempenho e energia em
fluxos de múltiplas requisições, entre eles o agêntico, em resumo estendido. Os
dois trabalhos se relacionam com a presente proposta pelo recorte: o primeiro
isola o custo de \textit{software} que envolve a computação do modelo, e o
segundo trata a sequência de requisições, e não a chamada isolada, como unidade
de análise. Diferem, contudo, no objeto, pois tratam da pilha de inferência e
não de uma camada da aplicação, situada acima dela.

Mais próximos do objeto deste trabalho estão os estudos que comparam formas de
orquestrar o modelo na aplicação. \citeonline{matinpour2026} avaliam dois
\textit{frameworks} de orquestração, LangChain e LlamaIndex, com oito modelos,
em tarefas de perguntas e respostas com recuperação aumentada e de raciocínio
matemático, medindo latência, acurácia, uso de recursos e consumo de energia.
Concluem que a escolha do \textit{framework} afeta o desempenho do sistema
independentemente do modelo: nas tarefas com recuperação, por exemplo, o
LlamaIndex teve latência de 3\% a 15\% menor e consumo de energia de 2\% a 4\%
maior. \citeonline{tripathy2026} avaliam quatro \textit{frameworks} agênticos
de correção de \textit{issues} (SWE-Agent, OpenHands, Mini SWE Agent e
AutoCodeRover) com dois modelos de linguagem pequenos, em \textit{hardware}
fixo e com 150 execuções por configuração, e concluem que a arquitetura do
\textit{framework} é o principal fator do consumo: o mais custoso consumiu, em
média, 9,4 vezes a energia do menos custoso, embora quase nenhuma execução
tenha resolvido a tarefa (4\% em uma configuração e 0\% nas demais), o que
indica que boa parte dessa energia foi gasta em ciclos de raciocínio
improdutivos. Em trabalho anterior, do qual o autor é coautor,
\citeonline{barros2026} compararam, por telemetria NVML (\textit{NVIDIA Management Library}, interface
de leitura de potência e energia da GPU; Seção~\ref{sec:energia}) em uma GPU
H100, o
processamento em um único \textit{prompt} com um debate entre agentes, em 100
tarefas de contexto longo, contabilizando economia de energia apenas quando a
configuração multiagente mantinha a qualidade da monolítica.

Os três trabalhos mostram que decisões de orquestração tomadas na aplicação
alteram o consumo de energia, o que sustenta a hipótese deste trabalho, e o
último fornece o protocolo de medição reaproveitado aqui. Diferenciam-se da
presente proposta pelo que comparam: \textit{frameworks} inteiros, que diferem
em muitas decisões ao mesmo tempo, em \citeonline{matinpour2026} e
\citeonline{tripathy2026}; topologias de agentes, em \citeonline{barros2026}.
Nenhum deles compara a mesma aplicação com e sem uma camada de integração,
mantendo fixa a orquestração, e nenhum avalia o esforço de trocar o provedor do
modelo.

Na literatura nacional, a busca na SBC OpenLib encontrou poucos trabalhos sobre
o tema. Além de \citeonline{barros2026}, \citeonline{nascimento2026} medem o
consumo de energia e a pegada de carbono de quatro arquiteturas de aprendizado
profundo em CPU local e em TPU na nuvem, com a biblioteca CodeCarbon, e
mostram que a escolha do \textit{hardware} e a localização da matriz elétrica
alteram esses valores. O trabalho compartilha a preocupação de tratar o
consumo como critério de decisão de projeto. A decisão estudada, porém, é de
infraestrutura, e não de arquitetura de \textit{software}.

% -----------------------------------------------------------------------------
\section{Síntese e posicionamento}
\label{sec:tr-sintese}

O Quadro~\ref{qua:trabalhos-relacionados} resume a comparação. As propostas
arquiteturais para integrar modelos de linguagem
\cite{liu2025,bucaioni2025,bandeira2026,morais2026,kandogan2025}, os relatos de
\textit{gateways} \cite{sheneman2026,rahaman2026} e o levantamento de
\textit{frameworks} \cite{zhu2026} discutem ou demonstram os benefícios de uma
camada de integração sem, segundo o que relatam, medi-los contra a integração
direta. Fora da IA,
\citeonline{mo2023} mediram o benefício e o custo de uma camada de abstração,
mas em desempenho e sobre serviços de nuvem, e \citeonline{xiao2025}
confrontaram energia e manutenibilidade de decisões arquiteturais, sem modelo
de linguagem e sem propor um critério para decidir entre as duas. Os estudos de energia em sistemas com LLMs, por sua vez, medem a
pilha de inferência ou comparam \textit{frameworks} e topologias inteiros
\cite{vellaisamy2026,ifath2026,matinpour2026,tripathy2026,barros2026}.

\begin{quadro}[h!]
  \centering
  \Caption{\label{qua:trabalhos-relacionados}Comparação entre os trabalhos
  relacionados e a proposta quanto ao objeto e às dimensões avaliadas}
  \UFCqua{}{
    \scriptsize
    \begin{tabularx}{\textwidth}{|>{\raggedright\arraybackslash}p{2.3cm}|>{\raggedright\arraybackslash}X|>{\raggedright\arraybackslash}p{1.9cm}|>{\raggedright\arraybackslash}p{1.4cm}|>{\raggedright\arraybackslash}p{1.6cm}|>{\raggedright\arraybackslash}p{1.2cm}|>{\raggedright\arraybackslash}p{1.4cm}|}
      \hline
      \textbf{Trabalho} & \textbf{Objeto} & \textbf{Tipo de avaliação} & \textbf{Esforço de troca me\-dido} & \textbf{Manu\-tenibi\-lidade me\-dida} & \textbf{Energia medida} & \textbf{Compara com e sem a decisão} \\
      \hline
      \citeonline{liu2025} & Padrões para agentes baseados em modelos de fundação & Catálogo derivado de revisão & Não & Qualitativo & Não & Não \\
      \hline
      \citeonline{bucaioni2025} & Arquitetura de referência para sistemas integrados a LLMs & Mapeamento a três sistemas & Não & Qualitativo & Não & Não \\
      \hline
      \citeonline{bandeira2026} & Arquitetura de referência para aplicações com o LLM como núcleo & Ilustração & Não & Não & Não & Não \\
      \hline
      \citeonline{morais2026} & Composição declarativa de agentes, independente do provedor & Estudo de viabilidade & Não (resumo) & Não (resumo) & Não & Não \\
      \hline
      \citeonline{kandogan2025} & Arquitetura para IA composta com registros de agentes e dados & Caso de uso ilustrativo & Não & Não & Não & Não \\
      \hline
      \citeonline{sheneman2026}; \citeonline{rahaman2026} & \textit{Gateways} de acesso a modelos de vários provedores & Relato de implantação & Não & Não & Não & Não \\
      \hline
      \citeonline{lin2026} & Comportamento de \textit{gateways} comerciais de terceiros & Medição (auditoria) & Não & Não & Não & Não \\
      \hline
      \citeonline{zhu2026} & \textit{Frameworks} e protocolos multiagente & \textit{Survey} & Não & Qualitativo & Não & Não \\
      \hline
      \citeonline{kaur2017} & Interoperabilidade e portabilidade em nuvens & Revisão & Não & Qualitativo & Não & Não \\
      \hline
      \citeonline{mo2023} & Biblioteca de abstração e APIs nativas em funções \textit{serverless} & Experimento (42 estudantes) & Sim & Sim & Não (desem\-penho) & Sim \\
      \hline
      \citeonline{poy2025} & Padrões de projeto, \textit{code smells} e refatorações & Mapeamento sistemático & Não & Não & Sim (estudos primários) & Sim (estudos primários) \\
      \hline
      \citeonline{feitosa2017} & Padrões de projeto GoF e alternativas sem o padrão & Experimento & Não & Não & Sim & Sim \\
      \hline
      \citeonline{xiao2025} & Táticas e padrões em microsserviços & Experimento & Não & Sim & Sim & Sim \\
      \hline
      \citeonline{jarvenpaa2024} & Táticas arquiteturais verdes para sistemas com aprendizado de máquina & Catálogo derivado de revisão & Não & Não & Não & Não \\
      \hline
      \citeonline{vellaisamy2026} & Sobrecarga do hospedeiro na inferência de LLMs & Caracterização & Não & Não & Não (latência) & Não \\
      \hline
      \citeonline{ifath2026} & Fluxos de múltiplas requisições a LLMs & Caracterização (resumo estendido) & Não & Não & Sim & Não \\
      \hline
      \citeonline{matinpour2026} & \textit{Frameworks} de orquestração (LangChain e LlamaIndex) & Experimento & Não & Não & Sim & Não (\textit{frameworks} inteiros) \\
      \hline
      \citeonline{tripathy2026} & \textit{Frameworks} agênticos com modelos pequenos & Experimento & Não & Não & Sim & Não (\textit{frameworks} inteiros) \\
      \hline
      \citeonline{barros2026} & Prompt único e debate multiagente & Experimento & Não & Não & Sim & Sim \\
      \hline
      \citeonline{nascimento2026} & \textit{Hardware} e matriz elétrica para modelos de aprendizado profundo & Experimento & Não & Não & Sim & Não \\
      \hline
      \textbf{Proposta deste trabalho} & \textbf{Camada de \harness{} entre aplicação \textit{web} e provedor de IA} & \textbf{Experimento} & \textbf{Sim} & \textbf{Sim} & \textbf{Sim} & \textbf{Sim} \\
      \hline
    \end{tabularx}
  }{
    \Fonte{Elaborado pelo autor}
  }
\end{quadro}

No quadro, \textit{Não (resumo)} indica que o resumo do trabalho não relata a
dimensão; o texto integral de \citeonline{morais2026} não foi lido. Só o
trabalho de \citeonline{mo2023} e a proposta medem o esforço de trocar a
abstração, e só a proposta o faz sobre a camada entre aplicação e provedor de
modelo.

Nenhum dos trabalhos encontrados reúne as três características da proposta:
ter como objeto a camada entre a aplicação e o provedor do modelo, medir o
esforço de trocar o provedor e medir a energia que a camada acrescenta,
comparando a mesma aplicação com e sem ela. É essa combinação que define a
contribuição do trabalho. Dos trabalhos analisados, ele herda o desenho com e
sem a camada e as métricas de código de \citeonline{mo2023}, sem herdar a
medida do esforço com participantes nem o uso de uma biblioteca pronta
(Seção~\ref{sec:tr-lockin}), a pergunta que confronta energia e
manutenibilidade de \citeonline{xiao2025}, as responsabilidades da camada do
catálogo de \citeonline{liu2025} e o protocolo de medição de
\citeonline{barros2026}. Os conceitos que sustentam essas
escolhas são apresentados no Capítulo~\ref{cap:fundamentacao}.
